\chapter{Drehimpuls, Spin und Statistik}

\section{Von der gewöhnlichen Rotation zum intrinsischen Spin}

Formal geht es um den Drehimpuls, eigentlich aber um den Spin. Der Spin ist eine der einfachsten und zugleich am wenigsten intuitiven Eigenschaften elementarer Teilchen; seine Struktur ergibt sich aus der Mathematik des Drehimpulses.

Der Drehimpuls ist ein Vektor. Seine Richtung liegt entlang der Rotationsachse. Entlang dieser Achse gibt es zwei mögliche Richtungen, daher braucht man eine Konvention, um zwischen ihnen zu wählen: Krümmt man die Finger der rechten Hand in Drehrichtung, so zeigt der Daumen in Richtung des Drehimpulsvektors. Das ist eine Definition, kein Satz.

Der Betrag hängt davon ab, wie viel Masse vorhanden ist, wie diese Masse verteilt ist, wie groß das System ist und wie hoch seine Winkelgeschwindigkeit ist. Massenverteilung und Größe gehen in das Trägheitsmoment ein. Ein großer und ein kleiner Körper können daher denselben Drehimpuls haben und sich dennoch mit sehr unterschiedlichen Geschwindigkeiten und Energien drehen.

Ein gewöhnlicher zusammengesetzter Körper kann zwei begrifflich verschiedene Arten von Drehimpuls besitzen. Bewegt sich eine Tasse auf einer Kreisbahn um eine äußere Achse, hat sie Drehimpuls, weil sich ihr Schwerpunkt um diese Achse bewegt – selbst wenn sich die Tasse nicht um eine durch sie selbst verlaufende Achse dreht. Das ist Bahndrehimpuls. Spin bezeichnet die Rotation um eine innere Achse.

Die klare Definition verwendet das Schwerpunktsystem. Man wechselt in das Bezugssystem, in dem der Schwerpunkt ruht, sodass der gewöhnliche Gesamtimpuls verschwindet. Jeder Drehimpuls, der in diesem System übrig bleibt, ist Spin-Drehimpuls. Der Körper muss sich im Inneren nicht in Ruhe befinden; nur sein Schwerpunkt ruht.

Bei einem Basketball ist diese Unterscheidung unproblematisch. Einen ruhenden Basketball und denselben, in Rotation versetzten Basketball erkennt man offensichtlich als denselben Körper in zwei verschiedenen Zuständen. Die Quantenmechanik verkompliziert dieses Bild, weil der Drehimpuls diskret ist. Es gibt keine kontinuierliche Folge erlaubter Drehimpulse, die einen Zustand mit Drehimpuls null mit dem ersten rotierenden Zustand verbindet. Man kann daher fragen, ob der erste rotierende Zustand denselben Körper in Bewegung oder einen ganz anderen Körper darstellt.

Die Formulierung ist teilweise konventionell, aber die physikalische Frage ist präzise: Wie viel Energie ist erforderlich, um den ersten nicht verschwindenden Rotationszustand zu erreichen? Ist die Energielücke klein, bleibt der angeregte Körper leicht wiederzuerkennen. Ist sie enorm, kann die Anregung den Körper radikal verändern. Versucht man, einem Atom zu viel Rotationsanregung zu geben, können seine Elektronen abgestoßen werden und das Atom zerstören, anstatt dasselbe Atom ruhig rotieren zu lassen.

Betrachten wir nun ein Elektron und lassen seinen bekannten intrinsischen Spin vorerst außer Acht. Kann man ein Elektron um eine innere Achse in Rotation versetzen? Ein mathematischer Punkt besitzt keine sichtbare innere Struktur, die rotieren könnte. Experimente bestimmen die Größe eines Elektrons jedoch nicht, indem sie ihm eine kleine starre Oberfläche geben, deren Drehung man beobachten könnte. Falls das Elektron eine extrem kleine, aber endliche Größe hat, könnte seine erste zusätzliche Rotationsanregung schlicht zu viel Energie erfordern, um erreichbar zu sein. Da die erforderliche Energie mit abnehmender Größe steigt, könnte sie astronomisch hoch sein, vielleicht sogar in der Größenordnung der Planck-Energie liegen. Diese Abschätzung ist ausdrücklich vorläufig.\footnote{Eine kurze Passage unmittelbar nach dieser Abschätzung ist nicht verständlich genug, um daraus eine weitere inhaltliche Rekonstruktion abzuleiten.}

Der Betrag des Spins wird bei einem elementaren Teilchen folglich als feste Eigenschaft behandelt. Elektronen gibt es nicht mit einer Reihe zunehmend größerer intrinsischer Spinbeträge. Hätte eine ansonsten elektronenähnliche Anregung einen anderen Gesamtspin, würde man sie als anderes Teilchen klassifizieren.

Der feste Betrag legt die Richtung nicht fest. Stellen wir uns ein Teilchen namens \emph{Spintron} vor, das teilweise durch den Betrag seines Drehimpulses charakterisiert ist. Zeigt sein Spin entlang einer Achse, so verlangt die Rotationsinvarianz, dass es dieselbe Teilchenart auch mit einem entlang einer anderen Achse ausgerichteten Spin geben kann. Keine räumliche Richtung ist grundsätzlich ausgezeichnet. Der Spinbetrag bleibt fest, während sich seine Orientierung ändern kann. Das Problem enthält also zwei scheinbar widersprüchliche Bestandteile: einen quantisierten Betrag und frei drehbare Richtungen. Die Drehimpulsalgebra wird zeigen, wie beides miteinander vereinbar ist.

\begin{classroomqa}
\lecturetimestamp{00:12:25}
\audiencequestion{Warum benötigt man Energie, um ein kleines Teilchen in Rotation zu versetzen?}
\lecturerresponse{Für einen Körper der Masse \(M\) und des charakteristischen Radius \(R\) gilt für das Trägheitsmoment näherungsweise
\begin{equation}
I\sim cMR^2,
\end{equation}
wobei der Zahlenfaktor \(c\) von der Form und der Massenverteilung abhängt. Bewegt sich Materie in der Nähe des äußeren Randes mit der Geschwindigkeit \(v\), so ist ihr gewöhnlicher Impuls von der Größenordnung \(Mv\), und der Drehimpuls ist Impuls mal Abstand:
\begin{equation}
L\sim MvR.
\end{equation}
Äquivalent gilt \(L=I\omega\) mit \(v\sim\omega R\). Die kinetische Rotationsenergie ist
\begin{equation}
E_{\mathrm{rot}}=\frac{L^2}{2I}.
\end{equation}
Bei festem \(M\) und festem \(L\) gilt
\begin{equation}
E_{\mathrm{rot}}\sim\frac{L^2}{2cMR^2}\propto\frac{1}{R^2}.
\end{equation}
Je kleiner der Körper ist, desto mehr Energie kostet ein gegebener Drehimpuls. Diese Abhängigkeit vom inversen Quadrat ist klassisch. Die zusätzliche quantenmechanische Tatsache besteht darin, dass sich der Drehimpuls in diskreten Schritten ändert, sodass der erste erlaubte Schritt nicht beliebig klein sein kann.}
\end{classroomqa}

Die nun verwendete Mathematik ist abstrakt und anfangs recht unintuitiv. Interessant ist, dass sich aus einer kleinen Zahl algebraischer Regeln experimentell überprüfbare Tatsachen ergeben.

\section{Bahndrehimpuls}

Betrachten wir zunächst ein Teilchen, das sich auf einer Kreisbahn um einen Ursprung bewegt. Der Betrag seines Drehimpulses ist gleich dem Impuls mal dem senkrechten Abstand vom Ursprung. In einer Ebene kann man eine Umlaufrichtung als positiv und die entgegengesetzte als negativ bezeichnen. Das Vorzeichen gibt die Richtung des Drehimpulsvektors relativ zur gewählten Achse an.

Das ist noch kein Spin. Das Teilchen hat Drehimpuls, weil es gewöhnlichen Impuls besitzt. Nun platzieren wir zwei Teilchen an gegenüberliegenden Enden eines Durchmessers und lassen sie sich im gleichen Drehsinn bewegen. Ihre gewöhnlichen Impulse heben sich auf, sodass der Schwerpunkt ruht; ihre Drehimpulse zeigen jedoch in dieselbe Richtung und addieren sich. Das System hat verschwindenden Gesamtimpuls und einen von null verschiedenen Gesamtdrehimpuls: ein klassisches Spinmodell.

Der Ortsvektor eines Teilchens relativ zum Ursprung sei
\begin{equation}
\vec r=(x,y,z)=(x_1,x_2,x_3),
\end{equation}
und sein Impuls
\begin{equation}
\vec p=(p_x,p_y,p_z)=(p_1,p_2,p_3).
\end{equation}
Der Drehimpuls ist das Vektorprodukt
\begin{equation}
\vec L=\vec r\times\vec p.
\end{equation}
Vertauscht man die Faktoren, kehrt sich der Vektor um: \(\vec p\times\vec r=-\vec r\times\vec p\). Die Standardkonvention lautet \(\vec r\times\vec p\); die Richtung ergibt sich aus der Rechte-Hand-Regel.\footnote{\raggedright Videobild aus der Vorlesung \nolinkurl{lecture_07_figure_02.png}, 00:18:49.\par}

\begin{figure}[h]
\centering
\includegraphics[width=0.72\textwidth]{lecture_07_figure_02.png}
\par
\begin{tikzpicture}[scale=1.0]
\draw[->] (-0.6,0) -- (4.8,0) node[right] {$x$};
\draw[->] (0,-0.6) -- (0,3.8) node[above] {$y$};
\coordinate (P) at (3.2,2.3);
\draw[->,thick] (0,0) -- (P) node[midway,above left] {$\vec r$};
\draw[dashed] (P) -- (3.2,0);
\draw[dashed] (P) -- (0,2.3);
\draw[->,thick] (P) -- ++(-1.0,1.25) node[above] {$\vec p$};
\node at (4.15,2.75) {$\vec L=\vec r\times\vec p$};
\end{tikzpicture}
\caption{Für ein Teilchen am Ort \(\vec r\) mit dem Impuls \(\vec p\) steht der Drehimpuls \(\vec L=\vec r\times\vec p\) senkrecht auf ihrer Ebene.}
\lectureframe{00:18:49}
\end{figure}

In kartesischen Komponenten gilt
\begin{equation}
\begin{aligned}
L_x&=yp_z-zp_y,\qquad
L_y=zp_x-xp_z,\\
L_z&=xp_y-yp_x.
\end{aligned}
\end{equation}
Es genügt, sich die erste Formel zu merken und die Bezeichnungen zyklisch zu vertauschen:
\[
x\longrightarrow y\longrightarrow z\longrightarrow x.
\]
Vertauscht man sowohl die Komponentenbezeichnung als auch die Bezeichnungen auf der rechten Seite zyklisch, erhält man die nächste Formel, ohne das relative Vorzeichen zu ändern. Diese Merkhilfe ist gerade deshalb nützlich, weil man das Vorzeichen leicht versehentlich umkehrt.

Für ein Punktteilchen ist \(\vec r\times\vec p\) der Bahndrehimpuls bezüglich des gewählten Ursprungs. Bei einem zusammengesetzten System addiert man die Drehimpulse aller Bestandteile als Vektoren:
\begin{equation}
\vec L_{\mathrm{tot}}=\sum_a \vec r_a\times\vec p_a.
\end{equation}
Wertet man diese Summe im Schwerpunktsystem aus, kann sie von null verschieden bleiben, obwohl
\[
\sum_a\vec p_a=0.
\]
Dieser verbleibende Drehimpuls ist der Spin des zusammengesetzten Systems.

\section{Quantenoperatoren und kanonische Kommutatoren}

In der Quantenmechanik werden beobachtbare Größen durch Operatoren dargestellt. Ort, Impuls und jede Drehimpulskomponente sind daher Operatoren. Ihre konkrete Realisierung als Differentialoperatoren wird hier nicht benötigt. Die Kommutationsrelationen reichen aus, um das Drehimpulsspektrum herzuleiten.

Für zwei Operatoren \(A\) und \(B\) gilt
\begin{equation}
[A,B]=AB-BA.
\end{equation}
Kommutieren zwei Observablen, können sie gleichzeitig festgelegt werden. Kommutieren sie nicht, verhindert die Quantenmechanik, dass beide zugleich scharf festgelegt sind.

Alle drei Ortskoordinaten können gemeinsam gemessen werden; daher kommutieren die Ortskomponenten:
\begin{equation}
[x_i,x_j]=0.
\end{equation}
Insbesondere kommutiert jeder Operator mit sich selbst. Es wäre sinnlos, wenn eine Observable mit einer weiteren Messung derselben Observable unverträglich wäre.

Auch die Impulskomponenten kommutieren:
\begin{equation}
[p_i,p_j]=0.
\end{equation}
Eine Ortskomponente kommutiert mit einer Impulskomponente in einer anderen Richtung, nicht aber mit der entsprechenden Impulskomponente. Die drei kanonischen Aussagen lassen sich gemeinsam schreiben als
\begin{equation}
[x_i,x_j]=0,\qquad
[p_i,p_j]=0,\qquad
[x_i,p_j]=i\hbar\,\delta_{ij}.
\end{equation}
Somit gilt zum Beispiel
\[
[x,p_y]=0,\qquad [x,p_x]=i\hbar,
\]
und bei umgekehrter Reihenfolge wechselt das Vorzeichen:
\begin{equation}
[p_j,x_i]=-i\hbar\,\delta_{ij}.
\end{equation}

Drehimpuls und Planck-Konstante haben dieselbe Dimension. Der Drehimpuls ist Abstand mal Impuls, während die Unschärferelation die Dimensionsform
\begin{equation}
\Delta x\,\Delta p\sim\hbar
\end{equation}
bis auf den bekannten Zahlenfaktor in der präzisen Ungleichung hat. Daher liegt es nahe, dass der Drehimpuls in Einheiten von \(\hbar\) quantisiert ist, auch wenn diese Quantisierung noch nicht hergeleitet wurde.

\section{Herleitung der Drehimpulsalgebra}

Als Erstes berechnen wir den Kommutator zweier Drehimpulskomponenten. Aus
\[
L_x=yp_z-zp_y,\qquad
L_y=zp_x-xp_z,
\]
folgt
\begin{align}
[L_x,L_y]
&=[yp_z-zp_y,\;zp_x-xp_z] \\
&=[yp_z,zp_x]-[yp_z,xp_z] \notag\\
&\quad{}-[zp_y,zp_x]+[zp_y,xp_z].
\end{align}

\begin{classroomqa}
\lecturetimestamp{00:31:41}
\audiencequestion{Kann man sich die Komponentenformeln merken, indem man die Symbole von links nach rechts zyklisch vertauscht, statt die drei Gleichungen von oben nach unten zyklisch durchzugehen?}
\lecturerresponse{Ja. Beide zyklischen Merkhilfen sind zulässig. Sie dienen lediglich dazu, die Komponenten in der richtigen zyklischen Reihenfolge zu bilden und dabei das Vorzeichen im Blick zu behalten.}
\end{classroomqa}

Einen Kommutator kann man verstehen, indem man fragt, ob sich die Faktoren in einem Produkt durch die Faktoren im anderen hindurchschieben lassen. In
\[
[yp_z,xp_z],
\]
kommutiert der Operator \(y\) sowohl mit \(x\) als auch mit \(p_z\), während \(p_z\) sowohl mit \(x\) als auch mit sich selbst kommutiert. Es gibt also kein Hindernis, und dieser Term verschwindet. Mit derselben Überlegung erhält man
\[
[zp_y,zp_x]=0.
\]
Nur zwei Produktpaare bleiben übrig.

Beim ersten ist das einzige Hindernis, dass \(p_z\) an \(z\) vorbeigeschoben werden muss:
\begin{align}
[yp_z,zp_x]
&=yp_zzp_x-zp_xyp_z \\
&=y(p_zz-zp_z)p_x \\
&=y[p_z,z]p_x \\
&=-i\hbar\,yp_x.
\end{align}
Die Reihenfolge im kanonischen Kommutator ist entscheidend:
\[
[p_z,z]=-[z,p_z]=-i\hbar.
\]

Beim anderen verbleibenden Paar ergeben die Minuszeichen in den Definitionen von \(L_x\) und \(L_y\) ein Plus. Das einzige Hindernis ist, dass \(z\) an \(p_z\) vorbeigeschoben werden muss:
\begin{align}
[zp_y,xp_z]
&=zp_yxp_z-xp_zzp_y \\
&=x(zp_z-p_zz)p_y \\
&=x[z,p_z]p_y \\
&=i\hbar\,xp_y.
\end{align}
Fasst man beide Beiträge zusammen, erhält man
\begin{align}
[L_x,L_y]
&=-i\hbar\,yp_x+i\hbar\,xp_y \\
&=i\hbar(xp_y-yp_x) \\
&=i\hbar L_z.
\end{align}

Das ist die entscheidende Vereinfachung. Der Kommutator zweier Drehimpulskomponenten ergibt die dritte Komponente; kein Orts- oder Impulsoperator bleibt übrig. Durch zyklische Vertauschung erhält man
\begin{equation}
\begin{aligned}
[L_x,L_y]&=i\hbar L_z,\qquad
[L_y,L_z]=i\hbar L_x,\\
[L_z,L_x]&=i\hbar L_y.
\end{aligned}
\end{equation}
Äquivalent dazu gilt
\begin{equation}
[L_i,L_j]=i\hbar\,\epsilon_{ijk}L_k.
\end{equation}

Die Algebra schließt sich auf den Drehimpulsoperatoren selbst. Nachdem diese Relationen hergeleitet sind, können die zu ihrer Herleitung verwendeten \(x_i\) und \(p_i\) beiseitegelassen werden. Dieselbe Algebra gilt, wenn der Gesamtdrehimpuls eine Summe über viele Bestandteile ist. Insbesondere gilt sie im Schwerpunktsystem, in dem der verbleibende Drehimpuls der Spin ist. Die Drehimpulskommutatoren lassen sich daher abstrakt verwenden, ohne den Spin zunächst als Bewegung kleiner Bestandteile zu konstruieren.

Die Wahl der \(x\)-, \(y\)- und \(z\)-Achsen ist beliebig. Dreht man zu einem anderen Satz zueinander senkrechter inertialer kartesischer Achsen, erfüllen die gestrichenen Komponenten dieselben Relationen:
\begin{equation}
[L'_x,L'_y]=i\hbar L'_z,
\end{equation}
sowie die entsprechenden zyklischen Relationen. Die Algebra ist rotationsinvariant, daher können physikalische Vorhersagen nicht davon abhängen, welche inertialen Achsen für die Rechnung gewählt wurden.

\begin{classroomqa}
\lecturetimestamp{00:39:09}
\audiencequestion{Ist der Spin-Drehimpuls sowohl unabhängig von Translationen als auch von Drehungen der Koordinatenachsen?}
\lecturerresponse{Ja. Der intrinsische Spin ist translationsunabhängig. Der Bahndrehimpuls bezüglich eines Ursprungs kann sich ändern, wenn man den Ursprung verschiebt; der im Schwerpunktsystem verbleibende Spin hängt jedoch nicht von dieser Verschiebung ab.}
\end{classroomqa}

Abstrakte Algebra kann endlos viele Identitäten hervorbringen, von denen die meisten physikalisch uninteressant sind. Ihr Wert zeigt sich, wenn eine Identität zu einer messbaren Tatsache wird oder etwas über die operationelle Bedeutung der Messung erkennen lässt. Bis hierher ist die Konstruktion noch hauptsächlich algebraisch. Im nächsten Schritt wird die diskrete Struktur des Spins herausgearbeitet.

\begin{classroomqa}
\lecturetimestamp{00:40:42}
\audiencequestion{Was geschieht, wenn sich die gewählten \(x\)-, \(y\)- und \(z\)-Achsen mit den Teilchen mitdrehen?}
\lecturerresponse{Die Betrachtung ist auf Inertialsysteme beschränkt. Ein rotierendes Koordinatensystem ist nichtinertial und führt zu Effekten wie Zentrifugal- und Corioliskräften. Der Drehimpuls ist gegenüber einem Wechsel in ein mitrotierendes Bezugssystem nicht invariant: Ein System, das sich in einem Inertialsystem sichtbar dreht, kann in einem mit ihm rotierenden System in Ruhe erscheinen. Die hier verwendete Rotationsinvarianz bezieht sich auf unterschiedlich orientierte inertiale kartesische Achsen, nicht auf zeitlich rotierende Achsen.}
\end{classroomqa}

\section{Leiteroperatoren und Quantisierungsachse}

Wir führen die beiden Kombinationen ein:
\begin{equation}
L_+=L_x+iL_y,
\qquad
L_-=L_x-iL_y.
\end{equation}
Sie ähneln den Erzeugungs- und Vernichtungsoperatoren des harmonischen Oszillators. Ihre Aufgabe ist es, eine gewählte Drehimpulskomponente zu erhöhen oder zu erniedrigen.

Die \(z\)-Achse ist nicht ausgezeichnet, aber eine Achse muss gewählt werden, wenn eine Komponente angegeben werden soll. Wir wählen \(z\) als Quantisierungsachse und schreiben
\begin{equation}
L_z\lvert m\rangle=m\hbar\,\lvert m\rangle.
\end{equation}
Damit ist \(m\) die gemessene \(z\)-Komponente in Einheiten von \(\hbar\). Während der Rechnung ist es oft praktisch, \(\hbar=1\) zu setzen; hier wird sie explizit beibehalten.

Der Kommutator \([L_+,L_-]\) ist proportional zu \(L_z\), wird für das unmittelbare Argument aber nicht benötigt. Zunächst kommt es auf den Kommutator jedes Leiteroperators mit \(L_z\) an. Aus
\[
[L_x,L_z]=-i\hbar L_y,
\qquad
[L_y,L_z]=i\hbar L_x,
\]
erhält man
\begin{align}
[L_+,L_z]
&=[L_x+iL_y,L_z] \\
&=[L_x,L_z]+i[L_y,L_z] \\
&=-i\hbar L_y+i(i\hbar L_x) \\
&=-\hbar(L_x+iL_y) \\
&=-\hbar L_+.
\end{align}
Ändert man das Vorzeichen des Terms \(iL_y\), ergibt sich
\begin{equation}
[L_-,L_z]=+\hbar L_-.
\end{equation}
Bei umgekehrter Reihenfolge in jedem Kommutator erhält man äquivalent
\begin{equation}
[L_z,L_+]=\hbar L_+,
\qquad
[L_z,L_-]=-\hbar L_-.
\end{equation}

Gleichzeitig kann nur eine Drehimpulskomponente scharf sein. Je zwei verschiedene Komponenten kommutieren nicht, sodass sie nicht gleichzeitig festgelegt werden können. Die gewählte scharfe Komponente legt die Quantisierungsachse fest; \(z\) zu wählen, ist eine Zweckmäßigkeit und keine physikalische Bevorzugung.

Nun wenden wir \(L_+\) auf einen Eigenzustand von \(L_z\) an. Der Kommutator lässt sich umschreiben als
\begin{equation}
L_zL_+=L_+L_z+\hbar L_+.
\end{equation}
Daher gilt
\begin{align}
L_z\bigl(L_+\lvert m\rangle\bigr)
&=\bigl(L_+L_z+\hbar L_+\bigr)\lvert m\rangle \\
&=m\hbar L_+\lvert m\rangle+\hbar L_+\lvert m\rangle \\
&=(m+1)\hbar\,L_+\lvert m\rangle.
\end{align}
Sofern er nicht verschwindet, ist \(L_+\lvert m\rangle\) ein Eigenzustand von \(L_z\), dessen Eigenwert um eine Einheit von \(\hbar\) gestiegen ist. Entsprechend gilt
\begin{equation}
L_z\bigl(L_-\lvert m\rangle\bigr)
=(m-1)\hbar\,L_-\lvert m\rangle.
\end{equation}
Somit erniedrigt \(L_-\) den Wert von \(m\) um eine Einheit.

Das ist bereits eine nichttriviale Schlussfolgerung. Welche Form das Spektrum von \(L_z\) auch immer hat, benachbarte Werte liegen um ganze Zahlen auseinander:
\begin{equation}
\ldots,\ m-2,\ m-1,\ m,\ m+1,\ m+2,\ \ldots
\end{equation}
Die Endpunkte sind noch nicht bestimmt.

\section{Endliche Multipletts und die Spinquantenzahl}

Eine bestimmte Teilchenart hat einen festen Drehimpulsbetrag. Ihre Drehimpulsrichtung kann geändert werden, doch eine unbegrenzte Vergrößerung des Betrags würde schließlich einen anderen Körper ergeben. Daher muss es für dieses Teilchen einen maximal möglichen Wert der \(z\)-Komponente geben.

Wie kann eine Leiter mit einem Erhöhungsoperator einen obersten Zustand haben? Der harmonische Oszillator liefert das Vorbild: Sein Vernichtungsoperator ergibt null, wenn der Grundzustand erreicht ist. Hier muss der Erhöhungsoperator null ergeben, wenn der oberste Zustand erreicht ist:
\begin{equation}
L_+\lvert m_{\max}\rangle=0.
\end{equation}
Die Aussage, dass \(L_+\) einen Eigenwert erhöht, gilt nur, wenn der dadurch erzeugte Vektor nicht verschwindet. Der Zustand mit höchstem Gewicht ist genau der Ausnahmezustand, an dem die Erhöhungsfolge endet.

Die Rotationsinvarianz legt den unteren Endpunkt fest. Kann der Drehimpuls maximal entlang \(+z\) zeigen, lässt sich derselbe physikalische Körper so drehen, dass er maximal entlang \(-z\) zeigt. Also gilt
\begin{equation}
m_{\min}=-m_{\max}.
\end{equation}
Das Spektrum ist symmetrisch um null und hat eine Einerschrittweite. Es gibt zwei mögliche Familien.

In der ersten Familie kommt null vor:
\begin{equation}
\begin{aligned}
m={}&-m_{\max},-m_{\max}+1,\ldots,-1,0,\\
&1,\ldots,m_{\max}-1,m_{\max}.
\end{aligned}
\end{equation}
Dann ist \(m_{\max}\) eine ganze Zahl, und die Zahl der Zustände beträgt
\begin{equation}
2m_{\max}+1.
\end{equation}
Der Maximalwert charakterisiert das Teilchen und heißt sein Spin. In der üblichen Schreibweise gilt
\begin{equation}
s\equiv m_{\max},
\qquad
\lvert s,m\rangle\equiv\lvert m\rangle
\end{equation}
innerhalb eines Multipletts mit festem Spin.\footnote{\raggedright Videobild aus der Vorlesung \nolinkurl{lecture_07_figure_03.png}, 01:01:08.\par}

\begin{figure}[htbp]
\centering
\includegraphics[width=0.70\textwidth]{lecture_07_figure_03.png}
\caption{Die Bedingung für den Zustand mit höchstem Gewicht beendet die Erhöhungsfolge; die erlaubten \(m\)-Werte bilden eine endliche Leiter mit Einerschrittweite.}
\lectureframe{01:01:08}
\end{figure}

Ein Teilchen mit Spin null hat nur \(m=0\). Ein Teilchen mit Spin eins hat
\[
m=1,0,-1,
\]
und ein Teilchen mit Spin zwei hat fünf Komponentenwerte. Das Higgs-Boson ist ein Beispiel für Spin null. Das Photon und das \(Z\)-Boson haben Spin eins. Ein Graviton hätte, falls es existiert, Spin zwei. Zusammengesetzte Körper können Spin drei und wesentlich höheren Spin haben: Atomkerne können das, ebenso hinreichend makroskopische rotierende Körper wie Basketballs und Galaxien, die beliebig große Drehimpulse tragen können. Was auf der üblichen Liste elementarer Teilchen fehlt, ist ein elementares Teilchen mit Spin drei.

Die andere symmetrische Familie mit Einerschrittweite enthält null nicht. Sie beginnt mit
\[
m=\frac12,-\frac12
\]
und setzt sich in Einerschritten fort:
\begin{equation}
\begin{aligned}
m={}&-m_{\max},-m_{\max}+1,\ldots,-\frac32,-\frac12,\\
&\frac12,\frac32,\ldots,m_{\max}.
\end{aligned}
\end{equation}
Hier ist \(m_{\max}\) eine halbe ganze Zahl. Das einfachste Mitglied hat Spin \(\tfrac12\), darauf folgen Spin \(\tfrac32\), Spin \(\tfrac52\) und so weiter. Drehimpulsmultipletts gehören daher entweder zur Familie mit ganzzahligem oder zur Familie mit halbzahligem Spin.

\begin{classroomqa}
\lecturetimestamp{01:03:46}
\audiencequestion{Hat auch ein Multiplett mit halbzahligem Spin einen Maximalwert?}
\lecturerresponse{Ja. Auch hier hat die Leiter einen Höchstwert. Für ein Teilchen mit Spin \(\tfrac32\) ist der Maximalwert \(m_{\max}=\tfrac32\); für ein Teilchen mit Spin \(\tfrac52\) ist er \(m_{\max}=\tfrac52\). Der einzige Unterschied besteht darin, dass der Maximalwert und alle anderen Werte im Multiplett halbe ganze Zahlen sind.}
\end{classroomqa}

\begin{classroomqa}
\lecturetimestamp{01:04:04}
\audiencequestion{Gibt es keine elementaren Teilchen mit Spin drei oder höher?}
\lecturerresponse{Im hier betrachteten relativistischen quantenmechanischen Rahmen gehören die Spins elementarer Teilchen auf der üblichen Liste zu den Werten \(0,\tfrac12,1,\tfrac32,\) und \(2\). Für zusammengesetzte Körper gilt diese Einschränkung nicht; sie können Spin drei oder einen noch größeren Spin haben.}
\end{classroomqa}

\begin{classroomqa}
\lecturetimestamp{01:04:50}
\audiencequestion{Stehen die negativen Einträge auf der Leiter für negativen Spin?}
\lecturerresponse{Nein. Ein negativer Wert von \(m\) ist ein Komponentenwert: Er besagt, dass der Drehimpuls relativ zur gewählten Achse in die entgegengesetzte Richtung zeigt. Der Spin des Teilchens ist der nichtnegative Maximalwert \(s=m_{\max}\).}
\end{classroomqa}

\begin{classroomqa}
\lecturetimestamp{01:05:10}
\audiencequestion{Was geschieht mit all den Orientierungen zwischen den diskreten Komponentenwerten?}
\lecturerresponse{Es gibt nicht nur zwei oder drei mögliche räumliche Orientierungen. Zwischenorientierungen werden durch quantenmechanische Überlagerungen der Zustände in der \(m\)-Basis dargestellt. Die folgende Konstruktion für Spin \(\tfrac12\) macht das explizit.}
\end{classroomqa}

\begin{classroomqa}
\lecturetimestamp{01:05:42}
\audiencequestion{Was ist mit schraubenförmiger Bewegung oder zirkular polarisierten Wellen?}
\lecturerresponse{Eine zirkular polarisierte elektromagnetische Welle ist aus zirkular polarisierten Photonen aufgebaut, und diese Photonen tragen Drehimpuls entlang der Ausbreitungsachse. Das Wort ``schraubenförmig'' verbindet Rotation mit Vorwärtsbewegung und ist daher keine passende Beschreibung für einen ruhend betrachteten Körper. Bei einem bewegten Teilchen ist die relevante Größe die Helizität: die Orientierung seines Spins relativ zu seiner Bewegungsrichtung.}
\end{classroomqa}

\begin{classroomqa}
\lecturetimestamp{01:07:10}
\audiencequestion{Ist ein bekanntes elementares Teilchen mit Spin \(\tfrac32\) bekannt?}
\lecturerresponse{Ein elementares Teilchen mit Spin \(\tfrac32\) ist bisher nicht direkt bestätigt worden. Die Supersymmetrie vermutet einen Partner des Gravitons mit Spin \(\tfrac32\), das Gravitino genannt wird; direkt nachgewiesen wurde es jedoch nicht.}
\end{classroomqa}

\begin{classroomqa}
\lecturetimestamp{01:08:01}
\audiencequestion{Wirken die Drehimpulsoperatoren nicht mehr auf die ursprüngliche Wellenfunktion beziehungsweise den ursprünglichen Ket-Raum?}
\lecturerresponse{Es war noch kein ursprünglicher Spinraum festgelegt worden. Die Algebra wird gerade dazu verwendet, diesen Raum zu bestimmen. Für festen Spin \(s\) hat er Basisvektoren, die durch die erlaubten Werte bezeichnet sind:
\[
\lvert s,m\rangle,\qquad m=-s,-s+1,\ldots,s.
\]
Dieser intrinsische Spinraum unterscheidet sich vom gewöhnlichen Raum räumlicher Wellenfunktionen, auf dem der Bahndrehimpuls als \(\mathbf r\times\mathbf p\) wirkt, obwohl beide dieselbe Drehimpulsalgebra realisieren.}
\end{classroomqa}

\section{Der Betrag des Drehimpulses}

\begin{classroomqa}
\lecturetimestamp{01:09:32}
\audiencequestion{Ist der Betrag des Drehimpulsvektors beobachtbar?}
\lecturerresponse{Ja. Am bequemsten verwendet man sein Quadrat:
\[
L^2=L_x^2+L_y^2+L_z^2.
\]
Sein Wert im Zustand mit höchstem Gewicht ist nicht einfach \(s^2\hbar^2\). Selbst wenn \(L_z\) scharf und maximal ist, können die nichtkommutierenden transversalen Komponenten nicht beide scharf null sein; sie tragen daher zusätzlich quantenmechanisch bei.}
\end{classroomqa}

Um diesen Beitrag zu berechnen, multiplizieren wir die Leiteroperatoren in der Reihenfolge, in der der Erhöhungsoperator zuerst auf ein Ket wirkt:
\begin{align}
L_-L_+
&=(L_x-iL_y)(L_x+iL_y) \\
&=L_x^2+L_y^2+iL_xL_y-iL_yL_x \\
&=L_x^2+L_y^2+i[L_x,L_y] \\
&=L_x^2+L_y^2-\hbar L_z.
\end{align}
Im letzten Schritt wird \([L_x,L_y]=i\hbar L_z\) verwendet. Daher gilt
\begin{equation}
L_x^2+L_y^2=L_-L_++\hbar L_z,
\end{equation}
und somit
\begin{equation}
L^2=L_-L_++L_z^2+\hbar L_z.
\end{equation}

Wenden wir diese Identität auf den Zustand mit höchstem Gewicht \(\lvert s,s\rangle=\lvert m_{\max}\rangle\) an:
\begin{align}
L^2\lvert s,s\rangle
&=\left(L_-L_++L_z^2+\hbar L_z\right)\lvert s,s\rangle \\
&=L_-L_+\lvert s,s\rangle
+s^2\hbar^2\lvert s,s\rangle
+s\hbar^2\lvert s,s\rangle.
\end{align}
Der erste Term verschwindet, weil \(L_+\lvert s,s\rangle=0\). Somit gilt
\begin{equation}
L^2\lvert s,s\rangle
=s(s+1)\hbar^2\lvert s,s\rangle.
\end{equation}
Das zusätzliche \(+s\) ist die quantenmechanische Korrektur, die mit den unvermeidlichen transversalen Fluktuationen verbunden ist. Klassisch hätte man nur \(s^2\hbar^2\) erwartet.

Die ersten Werte machen den Unterschied anschaulich:
\begin{align}
s=0 &: \qquad L^2=0,\\
s=1 &: \qquad L^2=2\hbar^2,\\
s=\frac12 &: \qquad L^2=\frac34\hbar^2.
\end{align}
Für \(s=1000\) beträgt der dimensionslose Faktor \(1000\times1001\); auf der Skala eines so großen Drehimpulses unterscheidet er sich nur wenig von \(1000^2\). Allgemein gilt
\begin{equation}
s(s+1)=s^2\left(1+\frac1s\right),
\end{equation}
sodass die relative quantenmechanische Korrektur für sehr große \(s\) vernachlässigbar wird. So ergibt sich das klassische Resultat \(s^2\).

Jeder \(m\)-Zustand in einem Multiplett mit festem Spin hat denselben Eigenwert von \(L^2\). Eine Änderung von \(m\) ändert die Komponente entlang der gewählten Achse, nicht aber den Gesamtbetrag. Das ist die quantenmechanische Entsprechung einer Drehung des Drehimpulsvektors, bei der seine Länge unverändert bleibt.

\begin{classroomqa}
\lecturetimestamp{01:16:43}
\audiencequestion{Können \(L^2\) und die \(z\)-Komponente gleichzeitig bekannt sein?}
\lecturerresponse{Ja. Das Gesamtquadrat kommutiert mit jeder Drehimpulskomponente:
\[
[L^2,L_i]=0,\qquad i=x,y,z.
\]
Daher können \(L^2\) und eine gewählte Komponente, etwa \(L_z\), gemeinsam festgelegt werden, obwohl zwei verschiedene Komponenten das nicht können. Jedes Elektron hat beispielsweise \(L^2=\tfrac34\hbar^2\), unabhängig von seiner Orientierung.}
\end{classroomqa}

\begin{classroomqa}
\lecturetimestamp{01:17:41}
\audiencequestion{Ist \(L^2\) bei Spin eins null oder nur bei Spin null?}
\lecturerresponse{Nur bei Spin null gilt \(L^2=0\). Für Spin eins ergibt sich
\[
L^2=1(1+1)\hbar^2=2\hbar^2,
\]
während man für Spin drei erhält
\[
L^2=3(3+1)\hbar^2=12\hbar^2.
\]}
\end{classroomqa}

\begin{classroomqa}
\lecturetimestamp{01:18:13}
\audiencequestion{Warum hat jede \(m\)-Stufe im Multiplett denselben Eigenwert von \(L^2\)?}
\lecturerresponse{Wir beginnen mit dem Zustand höchsten Gewichts und erniedrigen ihn. Da \(L^2\) mit jeder Komponente kommutiert, kommutiert er auch mit \(L_-\):
\[
[L^2,L_-]=0.
\]
Gilt
\[
L^2\lvert s,s\rangle=s(s+1)\hbar^2\lvert s,s\rangle,
\]
dann gilt
\begin{align}
L^2\bigl(L_-\lvert s,s\rangle\bigr)
&=L_-L^2\lvert s,s\rangle \\
&=s(s+1)\hbar^2\,L_-\lvert s,s\rangle.
\end{align}
Der nächste Zustand darunter hat somit denselben Eigenwert von \(L^2\). Wiederholt man dieses Argument, bleibt dieser Eigenwert entlang der gesamten Leiter erhalten.}
\end{classroomqa}

Die diskreten \(m\)-Werte bedeuten also nicht, dass es nur eine diskrete Menge räumlicher Orientierungen gibt. Die Rotationsinvarianz verlangt weiterhin Zustände, die gedrehten Richtungen entsprechen. Für Spin \(\tfrac12\) erscheinen diese Richtungen als Überlagerungen der beiden Basiszustände bezüglich \(L_z\).

\section{Spin-\texorpdfstring{\(\tfrac12\)}{1/2}-Zustände und Orientierung}

Die Rotationsinvarianz verlangt mehr als die diskreten Eigenzustände von \(L_z\). Ein Spin kann entlang jeder Achse zeigen, obwohl eine Messung seiner Komponente entlang einer gewählten Achse nur die erlaubten diskreten Werte ergibt. Die Richtungen zwischen den Eigenzuständen bezüglich der \(z\)-Achse werden durch Überlagerungen dargestellt.

Für ein Teilchen mit Spin \(\tfrac12\) betrachten wir nur die Spin-Freiheitsgrade. Ein Elektron hat außerdem Orbitalzustände und besitzt insgesamt also keineswegs nur zwei Zustände; allein sein Spinraum ist zweidimensional. Wähle die Basis bezüglich der \(z\)-Achse
\begin{equation}
L_z\lvert+\rangle=\frac{\hbar}{2}\lvert+\rangle,
\qquad
L_z\lvert-\rangle=-\frac{\hbar}{2}\lvert-\rangle.
\end{equation}
Ein allgemeiner Spinzustand lautet
\begin{equation}
\lvert\psi\rangle
=
\alpha\lvert+\rangle+\beta\lvert-\rangle,
\end{equation}
wobei \(\alpha\) und \(\beta\) komplexe Zahlen sind. Bei einer Messung von \(L_z\) erhält man mit der Wahrscheinlichkeit \(|\alpha|^2\) den Wert \(+\hbar/2\) und mit der Wahrscheinlichkeit \(|\beta|^2\) den Wert \(-\hbar/2\). Daher muss die Normierung gelten:
\begin{equation}
|\alpha|^2+|\beta|^2=1.
\end{equation}

Eine gemeinsame Phase hat keine physikalische Wirkung:
\begin{equation}
(\alpha,\beta)\longrightarrow
e^{i\chi}(\alpha,\beta).
\end{equation}
Die beiden komplexen Amplituden umfassen zunächst vier reelle Parameter. Durch die Normierung entfällt einer davon, und die irrelevante globale Phase beseitigt einen weiteren. Damit bleiben zwei physikalische reelle Parameter. Genau so viele braucht man, um eine Richtung durch ihren Polarwinkel und Azimutwinkel festzulegen. Eine Änderung von \(\alpha\) und \(\beta\) dreht die Spinorientierung.

Eine gemessene Komponente und ihr Erwartungswert sind zwei verschiedene Dinge. Der Erwartungswert ist ein Ensemblemittel. Man bereitet viele Elektronen exakt im selben Zustand vor, misst \(L_z\) einmal an jedem Elektron, addiert die Ergebnisse und teilt durch die Anzahl der Experimente. Jedes einzelne Messergebnis ist weiterhin entweder \(+\hbar/2\) oder \(-\hbar/2\), doch der Mittelwert kann jeden Wert zwischen diesen beiden annehmen:
\begin{equation}
\langle L_z\rangle
=
\frac{\hbar}{2}\bigl(|\alpha|^2-|\beta|^2\bigr).
\end{equation}

Gleiche Amplituden für Spin-up und Spin-down entsprechen einem Spin, der entlang einer der Orientierungen der \(x\)-Achse ausgerichtet ist. Normierte Zustände sind\footnote{Die Verhältnisse der Koeffizienten legen die Richtungen fest. Die Faktoren \(1/\sqrt2\) sorgen für die durch \( |\alpha|^2+|\beta|^2=1\) erforderliche Normierung. Die relativen Phasen \(+i\) und \(-i\) entsprechen entgegengesetzten Orientierungen entlang der \(y\)-Achse; welches Vorzeichen der positiven \(y\)-Richtung zugeordnet wird, hängt von der Achsenkonvention ab.}
\begin{equation}
\lvert x;+\rangle
=
\frac{\lvert+\rangle+\lvert-\rangle}{\sqrt2},
\qquad
\lvert x;-\rangle
=
\frac{\lvert+\rangle-\lvert-\rangle}{\sqrt2}.
\end{equation}
Für jeden der beiden Zustände liefert eine Messung entlang \(z\) die beiden möglichen Ergebnisse mit gleicher Wahrscheinlichkeit. Daher gilt
\begin{equation}
\langle L_z\rangle
=
\frac12\left(\frac{\hbar}{2}\right)
+
\frac12\left(-\frac{\hbar}{2}\right)
=0.
\end{equation}
Die relativen Vorzeichen plus und minus entsprechen entgegengesetzten Richtungen entlang der \(x\)-Achse. Entsprechend ergeben
\begin{equation}
\frac{\lvert+\rangle+i\lvert-\rangle}{\sqrt2},
\qquad
\frac{\lvert+\rangle-i\lvert-\rangle}{\sqrt2}
\end{equation}
die beiden entgegengesetzten Orientierungen entlang der \(y\)-Achse. Allgemeine normierte Werte von \(\alpha\) und \(\beta\) beschreiben beliebige Achsen. Zu sagen, ein Spin zeige entlang einer bestimmten Achse, bedeutet, dass eine Messung der Komponente entlang dieser Achse ein eindeutiges Ergebnis liefert. Bezüglich einer anderen Achse ist derselbe Zustand im Allgemeinen eine Superposition.

Die einzelnen Messergebnisse bleiben quantisiert. Stetig variieren die Ensemblemittelwerte; sie bilden die vertraute Geometrie eines klassischen Vektors nach.

\section{Spin, Statistik und zusammengesetzte Teilchen}

Masse, elektrische Ladung und Spin gehören zu den wichtigsten Eigenschaften, mit denen Teilchen charakterisiert werden. Der Spin steht außerdem mit der Statistik in Zusammenhang. Die Spin-Statistik-Beziehung ist ein Satz der relativistischen Quantenmechanik, wird hier aber nicht bewiesen:
\begin{equation}
\begin{aligned}
\text{half-integer spin} &\quad\Longleftrightarrow\quad \text{fermion},\\
\text{integer spin} &\quad\Longleftrightarrow\quad \text{boson}.
\end{aligned}
\end{equation}
Zwei identische Fermionen können nicht denselben vollständigen Quantenzustand einnehmen. Identische Bosonen können denselben Zustand einnehmen; große Ansammlungen von ihnen in einem Zustand können sich kollektiv wie eine klassische Welle verhalten. Teilchen mit Spin \(\tfrac12,\tfrac32,\tfrac52,\ldots\) sind also Fermionen, während Teilchen mit Spin \(0,1,2,\ldots\) Bosonen sind.

Die Regel lässt sich unmittelbar auf zusammengesetzte Teilchen übertragen. Die Addition zweier halbzahliger Drehimpulse ergibt einen ganzzahligen Drehimpuls. Allgemeiner bilden eine gerade Anzahl fermionischer Bestandteile ein zusammengesetztes Boson, eine ungerade Anzahl hingegen ein zusammengesetztes Fermion. Die Addition bosonischer Bestandteile ändert diese Paritätsregel nicht.

Gewöhnlicher Wasserstoff besteht aus einem Proton mit Spin \(\tfrac12\) und einem Elektron mit Spin \(\tfrac12\). Er enthält somit eine gerade Anzahl von Fermionen und ist als Ganzes bosonisch. Ein Deuteron besteht aus einem Proton und einem Neutron, die beide Fermionen sind, und ist ebenfalls ein Boson. Ein Deuteriumatom enthält zusätzlich zum Deuteron ein Elektron. Auf dieser Beschreibungsebene besteht es somit aus drei fermionischen Bestandteilen und ist daher ein Fermion.

Teilchen und Antiteilchen haben dieselbe Masse und denselben Spin, aber entgegengesetzte elektrische Ladungen. Positronium besteht aus einem Elektron und einem Positron, die aneinander gebunden sind. Es enthält zwei Fermionen und ist daher ein Boson.

\begin{classroomqa}
\lecturetimestamp{01:33:08}
\audiencequestion{Gehen Bosonen durch Dinge hindurch, und hängt Supraleitung damit zusammen, dass gepaarte Elektronen sich wie Bosonen verhalten?}
\lecturerresponse{Entscheidend ist nicht einfach die Fähigkeit, Materie zu durchdringen, sondern die Fähigkeit von Bosonen, denselben Zustand einzunehmen. Suprafluides Helium ist das einfachere Beispiel. Ein gewöhnliches Helium-\(4\)-Atom enthält zwei Protonen, zwei Neutronen und zwei Elektronen, also eine gerade Anzahl fermionischer Bestandteile. Das Atom ist daher ein Boson. Viele Heliumatome können denselben Quantenzustand einnehmen und sich kollektiv bewegen. Eine Supraflüssigkeit lässt sich dann näherungsweise als klassische Welle auffassen, die entsteht, wenn viele Atome in denselben Zustand gelangen. Diese Vorstellung ist nicht exakt, trifft aber den hier entscheidenden Punkt. Ein Elektron in einem Metall ist ein Fermion und kann keinen vollständigen Zustand mit einem anderen Elektron teilen. Paare aus Elektronen sind jedoch zusammengesetzte Bosonen; darin besteht der Zusammenhang mit der Supraleitung. Das wirft ein Paradox auf: Wenn die einzelnen Elektronen nicht denselben Zustand einnehmen können, wie können dann mehrere zusammengesetzte Paare denselben Zustand einnehmen? Zur Auflösung ist eine genauere Beschreibung zusammengesetzter Zustände erforderlich; darauf kommen wir später zurück.}
\end{classroomqa}

\begin{classroomqa}
\lecturetimestamp{01:35:46}
\audiencequestion{Da Protonen und Neutronen zusammengesetzt sind, warum haben sie dann halbzahligen Spin?}
\lecturerresponse{Beide bestehen aus einer ungeraden Anzahl von Quarks mit Spin \(\tfrac12\). Drei fermionische Bestandteile ergeben ein zusammengesetztes Fermion. Daher können Protonen und Neutronen halbzahligen Spin haben.}
\end{classroomqa}

\section{Pauli-Ausschlussprinzip und Atomorbitale}

Das Periodensystem liefert den unmittelbaren Anlass, das Pauli-Ausschlussprinzip einzuführen. Die Regel besagt nicht, dass zwei Elektronen niemals dieselbe Orbitalbewegung haben können. Sie besagt, dass zwei Elektronen nicht denselben vollständigen Quantenzustand einnehmen können.

Dieser Unterschied ist selbst bei den einfachsten Atomen wichtig. Man geht von Wasserstoff aus, erhöht die Kernladung und fügt ein zweites Elektron hinzu. Die beiden Elektronen können dasselbe räumliche Orbital einnehmen, ihre Spin-Zustände müssen sich jedoch unterscheiden. Bezüglich der gewählten Achse kann das eine Elektron Spin-up und das andere Spin-down haben. Nehmen zwei Elektronen dasselbe Orbital ein, bilden ihre Spins zusammen den Spin-Singulettzustand, sodass sich ihr Gesamtspin aufhebt. Ein einzelnes räumliches Orbital kann höchstens zwei Elektronen aufnehmen, denn es gibt nur zwei unabhängige Spin-Zustände.

\begin{classroomqa}
\lecturetimestamp{01:37:53}
\audiencequestion{Kann man sich die entgegengesetzten Spins zweier Elektronen in einem Orbital wie ineinandergreifende Zahnräder vorstellen?}
\lecturerresponse{Als grobe Vorstellung: ja. Ineinandergreifende Zahnräder drehen sich in entgegengesetzte Richtungen. Das ist lediglich ein Vergleich für die entgegengesetzten Spinrichtungen und kein wörtliches mechanisches Modell der Elektronen.}
\end{classroomqa}

Das Wort ``Orbital'' kann auf zwei verschiedenen Beschreibungsebenen verwendet werden. Ein Zustand der Orbitalbewegung im weiteren Sinn ist der Zustand, den man erhält, wenn man die Schrödingergleichung ohne Berücksichtigung des Spins löst. Er enthält mehr Informationen als nur die Angabe, welche Schale das Elektron besetzt.

\begin{classroomqa}
\lecturetimestamp{01:38:09}
\audiencequestion{Bedeutet ein Orbital mehr als nur die Angabe, welche Schale das Elektron besetzt?}
\lecturerresponse{Ja. Mit Orbitalbewegung ist hier der Schrödinger-Zustand des Elektrons gemeint, bei dessen Bestimmung der Spin unberücksichtigt bleibt. Dazu gehört die Hauptquantenzahl \(n\), die radiale Informationen enthält und ungefähr den Abstand des Elektrons vom Kern festlegt, ebenso wie die Quantenzahlen des Bahndrehimpulses. In der spektroskopischen Notation der Chemie entsprechen \(S\), \(P\), \(D\) und \(F\) jeweils dem Bahndrehimpuls \(l=0,1,2,3\).}
\end{classroomqa}

Die Hauptquantenzahl und der Bahndrehimpuls beschreiben also verschiedene Aspekte der Bewegung. Die Bezeichnungen \(S,P,D,F\) beziehen sich auf den Wert von \(l\), während \(n\) gesonderte radiale Informationen liefert. Paulis Regel besagt, dass keine zwei Elektronen in sämtlichen Quantenzahlen übereinstimmen dürfen, einschließlich der Spinquantenzahl.

\begin{classroomqa}
\lecturetimestamp{01:40:42}
\audiencequestion{Lässt sich mit den Quantenzahlen bestimmen, wie viele Elektronen sich in jeder Schale eines Atoms befinden können?}
\lecturerresponse{Ja. Sie legen die maximal mögliche Anzahl fest, aber nicht, wie viele Elektronen ein bestimmtes Atom tatsächlich in dieser Schale hat. Auf die allgemeine Frage, wie man die Schalenbesetzung zählt, gehen wir hier nicht weiter ein.}
\end{classroomqa}

Für einen festen Bahndrehimpuls \(l\) gibt es für die Quantenzahl der Komponente
\begin{equation}
2l+1
\end{equation}
mögliche Werte. Das ist dieselbe Zählung für Drehimpulse, die wir bereits beim Spin verwendet haben. Ein \(S\)-Zustand hat \(l=0\) und somit einen Orbitalzustand. Ein \(P\)-Zustand hat \(l=1\) und somit drei Orbitalzustände. Ein \(D\)-Zustand hat \(l=2\) und somit fünf.

Der Spin verdoppelt jede Anzahl von Orbitalzuständen, da jedes räumliche Orbital je ein Elektron in den beiden entgegengesetzten Spin-Zuständen aufnehmen kann.

\begin{classroomqa}
\lecturetimestamp{01:41:54}
\audiencequestion{Kann der \(S\)-Zustand zwei Elektronen aufnehmen, und können die drei \(P\)-Zustände zusammen sechs weitere aufnehmen?}
\lecturerresponse{Ja. Der einzelne \(S\)-Orbitalzustand kann höchstens zwei Elektronen aufnehmen. Die drei \(P\)-Orbitalzustände nehmen jeweils zwei Elektronen auf, also sechs weitere. Die hier betrachteten \(S\)- und \(P\)-Zustände können somit insgesamt \(2+6=8\) Elektronen aufnehmen.}
\end{classroomqa}

Dieselbe Drehimpulsstruktur, die die Atomorbitale ordnet, tritt mit einigen Unterschieden erneut auf, wenn sich Quarks in gebundenen Systemen bewegen und umeinander kreisen.